\chapter{Knee Injury}

\section{Definition}
Knee injuries involve damage to the bones, ligaments, tendons, cartilage, or menisci of the knee joint. They are common in athletes but can occur in anyone. Severity ranges from minor sprains to complex multi-ligament injuries.

\section{Pathophysiology}
The knee is a hinge joint stabilized by four major ligaments (ACL, PCL, MCL, LCL), two menisci (medial and lateral), and surrounding muscles and tendons. Injury mechanisms include twisting (non-contact ACL tear), direct impact (MCL tear, fracture), or overuse (patellar tendinitis, bursitis). Damage causes pain, swelling, instability, and impaired function.

\section{Etiology}
\begin{itemize}
\item Sports injuries: pivoting, cutting, landing from a jump, direct blows
\item ACL tear: non-contact twisting injury common in soccer, basketball, skiing
\item MCL tear: valgus stress (blow to the outer knee)
\item Meniscus tear: twisting with the foot planted
\item Patellar dislocation: twisting or direct blow
\item Patellofemoral pain syndrome: overuse, muscle imbalance
\item Fractures: patella, femoral condyles, tibial plateau
\item Overuse: runner's knee, jumper's knee
\item Osteoarthritis: degenerative changes
\item Direct trauma: falls, motor vehicle accidents
\end{itemize}

\section{Indications for Evaluation}
\begin{itemize}
\item Knee pain with weight-bearing or range of motion
\item Swelling (effusion) within hours of injury
\item Bruising around the knee
\item Instability or giving way of the knee
\item Locking or catching sensation (meniscus tear)
\item Inability to fully extend or flex the knee
\item Tenderness over specific structures (joint line, ligaments)
\item Audible pop at the time of injury (classic ACL tear)
\end{itemize}

\section{Related Symptoms}
\begin{itemize}
\item Knee swelling
\item Knee locking
\item Knee instability
\item Leg pain
\item Joint stiffness
\item Osteoarthritis (post-traumatic)
\end{itemize}

\section{Self-Care/Home Management}
\begin{itemize}
\item Rest, Ice, Compression, Elevation (RICE protocol)
\item Avoid weight-bearing on the injured knee
\item Use crutches for support
\item Take OTC pain relievers (ibuprofen, acetaminophen)
\item Apply ice for 20 minutes every 2-3 hours
\item Use an elastic bandage for compression
\item Begin gentle range-of-motion exercises once pain subsides
\end{itemize}

\section{Diagnostic Approach}
\begin{itemize}
\item History of injury mechanism and symptoms
\item Physical examination: inspection, palpation, range of motion, stability tests (Lachman, drawer, McMurray)
\item X-ray (weight-bearing) to rule out fracture
\item MRI for soft tissue evaluation (ligaments, menisci, cartilage)
\item CT scan for complex fractures
\item Diagnostic arthroscopy in selected cases
\end{itemize}

\section{Treatment/Management Overview}
\begin{itemize}
\item Conservative: RICE protocol, physical therapy, activity modification
\item Bracing for ligament injuries
\item Physical therapy for muscle strengthening and proprioception
\item Arthroscopic surgery for meniscus repair or partial meniscectomy
\item ACL reconstruction for active individuals
\item MCL healing with bracing (most heal non-operatively)
\item Knee replacement for severe arthritis
\end{itemize}
